\section{DenseShift: power-of-two weights}
\label{sec:ds}

The integer formats of \cref{sec:formats} kept one type for A, B, C and the sum, so they measure hardware cost but are
not what a quantised network computes. This section implements DenseShift~\cite{denseshift}, a low-bit
power-of-two weight format, end to end: training in PyTorch, a kernel with separate operand types and a shift
instead of a multiplier, builds on the card, and a small network whose layers run on the accelerator and give the
same predictions as the CPU. The work is on the git branch \code{denseshift}; the dated record is in
\file{docs/notes3/}.

\subsection{The method}
\label{sec:ds:method}
An \(n\)-bit DenseShift weight is \(w = s\cdot 2^{p}\) with a sign \(s = \pm 1\) and an exponent
\(p \in \{0,\dots,2^{n-1}-1\}\). Unlike earlier shift networks it has no zero, so all \(2^n\) codes carry a value:
2 bits give \(\{\pm1,\pm2\}\), 3 bits \(\{\pm1,\pm2,\pm4,\pm8\}\). The paper's three claims for hardware and training:
\begin{itemize}
  \item \textbf{No multiplier, no zero test.} With integer activations a MAC is \(c \mathrel{+}= \pm(x \ll p)\); without
        a zero weight the branch that shift networks need for \(w = 0\) disappears (Fig.~3 of the paper).
  \item \textbf{Float activations.} \(x\cdot(\pm 2^p)\) for a float \(x\) is an integer addition of \(p\) to the exponent
        field plus a sign flip (Eq.~2 of the paper), so even float networks need no multiplier.
  \item \textbf{Training.} The sign-scale decomposition writes \(w = (2H(w_s) - 1)\,2^{S_T}\) with
        \(S_0 = 0,\ S_t = H(w_t)(S_{t-1}+1)\), \(t = 1\dots T\), \(T = 2^{n-1}-1\) and the Heaviside step \(H\); each
        weight is trained through \(T+1\) float parameters with a straight-through estimator~\cite{s3shift}. The
        gradient to \(w_s\) is rescaled by \(\sqrt{S_T+1}\) instead of \(2^{S_T}\) (Eq.~8), and all parameters start
        from \(\mathcal N(0, 10^{-6})\) (``low-variance'' init).
\end{itemize}
The paper leaves some details to its released code, which was read for them: the forward value is \(s\,2^{S}\) and the
backward goes through \(s\sqrt{S+1}\); \(H(0) = 0\), \(\mathrm{sign}(0) = +1\); the first convolution and the final
fully connected layer stay float; there is no learned weight scale (BatchNorm after every convolution absorbs it).

\subsection{Training in PyTorch}
\label{sec:ds:train}
PyTorch comes from the system Python (torch 2.15 with CUDA on an RTX 5060 Ti); the flow's virtual environment is left
as it is, and the two exchange files. \file{ml/denseshift/denseshift.py} implements DenseShift linear and
convolution layers for any \(n \ge 2\). A self-test checks that only \(\pm2^p\) values occur, that every code is
reached, that the exported hardware codes decode to the trained weights, and that the 3-bit layer equals a
transcription of the official layer in values \emph{and} gradients. ResNet-20 on CIFAR-10 is a small-scale check of
the accuracy claim (the paper reports ImageNet), trained as in the official code: SGD 0.1, momentum 0.9, weight decay
\(10^{-4}\), cosine schedule, 200 epochs (\cref{tab:ds-cifar}).

\begin{table}[ht]
\centering\small
\caption{ResNet-20 on CIFAR-10, test accuracy after 200 epochs (seed 0; first convolution and FC layer float).
  Exported codes decoded into ordinary weights give the same accuracy in every run. Paper, ImageNet ResNet-18: float
  69.6, 2-bit 68.90, 3-bit 70.62, 4-bit 70.94\,\%.}
\label{tab:ds-cifar}
\begin{tabular}{@{}lrrl@{}}
\toprule
Weights & Accuracy & vs.\ float & Share of weights at \(|w| = 1, 2, 4, \dots\) \\
\midrule
float32 & 91.68\,\% & & \\
DenseShift 2-bit & 90.84\,\% & \(-0.84\) & 0.49 / 0.51 \\
DenseShift 3-bit & 91.40\,\% & \(-0.28\) & 0.52 / 0.14 / 0.12 / 0.21 \\
DenseShift 4-bit & 90.80\,\% & \(-0.88\) & 0.33 / 0.12 / 0.11 / 0.10 / 0.10 / 0.09 / 0.07 / 0.07 \\
\midrule
3-bit, Eq.~7 gradient & 91.36\,\% & \(-0.32\) & 0.76 / 0.18 / 0.03 / 0.03 \\
3-bit, Kaiming init & 90.50\,\% & \(-1.18\) & 0.45 / 0.20 / 0.15 / 0.20 \\
\bottomrule
\end{tabular}
\end{table}

3-bit DenseShift comes within 0.3 points of float. Without the rescaled gradient most exponents stay at zero (76\,\% of
the weights at \(\pm1\)), as the paper predicts, although the accuracy is the same here; the low-variance
initialisation is worth 0.9 points even when training from scratch. A second network, a multi-layer perceptron on
MNIST with two DenseShift hidden layers (\(784 \to 1024 \to 1024 \to 10\)), BatchNorm and unsigned 8-bit activations
learned with quantisation-aware training, reaches 98.65\,\% with float weights and 98.46/98.47/98.37\,\% with
2/3/4-bit DenseShift weights. It is the network that runs on the card (\cref{sec:ds:e2e}).

\subsection{A kernel with separate A, B and C types}
\label{sec:ds:kernel}
DenseShift needs what the single-type kernel cannot do: activations, weight codes and sums of different widths, and a
shift in place of the multiplier. Because the HLS hash covers the whole kernel tree, the change went into a copy,
\file{kernels/gemm\_mixed/}, so that every earlier build keeps its id; the flow picks the tree from the
configuration (\code{kernel.b\_type}, \code{b\_bits}, \code{accum\_bits}, hashed only when set).
\begin{itemize}
  \item \code{DataA\_t}, \code{DataB\_t}, \code{DataC\_t} as \code{ap\_[u]int} types (A also \code{float}); B and C
        get their own memory widths. 3-bit codes are stored in 4-bit lanes (128 per 512-bit word).
  \item The MAC is a function chosen in \file{Config.h}: multiply (\(c + a\cdot b\) in exact widths) or DenseShift,
        \(c + (s\,?\,-(a\ll e) : (a \ll e))\). Negating the narrow shifted term before one C-wide adder is 12\,\%
        smaller (HLS estimate) than an adder, a subtractor and a multiplexer.
  \item With 32-bit sums a PE's row of C (\(P_M\) lanes) is wider than a memory word, so the C path either splits a
        row over several words, passes it through or gathers several rows, all in one pipelined loop.
  \item A testbench draws codes over their full range and compares with an exact reference; 16 cases (every MAC,
        signed and unsigned activations, 16- and 32-bit sums with wrap-around, all three C conversions, whole and
        partial tiles, float activations) pass at two sizes each. The host runner gained \code{--b-dtype},
        \code{--c-dtype} and file input/output of codes, with a model checked against \code{ap\_int} arithmetic.
\end{itemize}

\paragraph{C buffers in BRAM.} The first HLS run of int8\,\(\times\)\,3-bit\,\(\to\)\,int32 at PE 32\(\times\)32
estimated 2002 BRAM18, more than the device has. Each PE's C buffer is 1024 bits wide and 512 deep. The utilisation of
an earlier int8 build showed why it is expensive: HLS's default two-port RAM becomes true-dual-port blocks with 36-bit
ports, which leave half of every BRAM36 empty at depth 512. The buffer needs only one write and one read port; bound as
a simple dual-port RAM (\code{BIND\_STORAGE ram\_s2p}) it uses 72-bit blocks, and the built design has 14 RAMB36 +
1 RAMB18 per PE instead of about 29. Even so, with 32-bit sums the C buffers are what limits the array size.

\subsection{Shift against multiply on the card}
\label{sec:ds:card}
\Cref{tab:ds-same} compares the two MACs at one array with the same activation and sum widths. Both run at the same
speed, as they must (same array, same schedule), and both are exact on the card.

\begin{table}[ht]
\centering\small
\caption{int8 activations and int32 sums, PE 32\(\times\)32, tile 512\(\times\)1024, transposed A, 200\,MHz, both
  meeting timing. HLS: estimates; Vivado: whole design after implementation; GOP/s at \(8192^3\), verified exactly.}
\label{tab:ds-same}
\begin{tabular}{@{}lrrrrrrr@{}}
\toprule
& \multicolumn{2}{c}{HLS} & \multicolumn{4}{c}{Vivado} & \\
\cmidrule(lr){2-3}\cmidrule(lr){4-7}
B (weights), MAC & DSP & LUT & DSP & LUT & FF & BRAM & GOP/s \\
\midrule
int8, multiply & 1064 & 195.8k & 1064 & 167.5k & 259.4k & 638.5 & 375.45 \\
DenseShift 3-bit, shift & 39 & 290.6k & 40 & 161.6k & 240.2k & 636.5 & 375.45 \\
\bottomrule
\end{tabular}
\end{table}

With a 32-bit sum HLS maps each 8-bit multiply and its add into one DSP48E1 (the single-type int8 kernel, whose
product was cut to 8 bits, had its multipliers in LUTs). The shift MAC removes all 1024 of them. HLS estimated that it
would cost 95k more LUTs instead; Vivado's implementation uses 5.9k \emph{fewer} LUTs and 19k fewer flip-flops (per
PE: 3433 against 4026 LUTs). On this FPGA DenseShift therefore frees the DSPs at no cost in logic. The array cannot
simply grow into them, because the 32-bit C buffers already use most of the BRAM (\cref{sec:ds:kernel}).

\subsection{End to end: an MNIST network on the card}
\label{sec:ds:e2e}
\file{ml/denseshift/mlp\_mnist.py} exports the MLP as an integer pipeline: activation codes \(\times\) DenseShift codes
give exact int32 sums; BatchNorm and the activation scale fold into a per-channel \(a, b\), and the next layer's codes
are \(\mathrm{clip}(\mathrm{rint}((c\,a + b)/s), 0, q_{\max})\). Evaluated on the CPU this pipeline has the same test
accuracy as the PyTorch model. On the card, both hidden layers run as GEMMs over all 10\,000 test images, once with
uint8 activation codes on a uint8\,\(\times\)\,3-bit build (same array; its default-preset build missed timing by
6\,ps on the PCIe user clock, the \code{explore} preset closed it) and once with codes 0..127 on the signed int8 build
of \cref{tab:ds-same}:
\begin{itemize}
  \item layer 1, \(10048\times832\times1024\) (padded): 52.6\,ms, 325.6\,GOP/s; layer 2,
        \(10048\times1024\times1024\): 63.2\,ms, 333.7\,GOP/s;
  \item in both variants 0 of 10\,240\,000 sums differ from the CPU in either layer; the test accuracy computed from
        the card's sums is 98.47\,\% (uint8) and 98.48\,\% (0..127), equal to PyTorch, and every prediction equals
        the CPU pipeline's.
\end{itemize}

\subsection{Float activations without a multiplier}
\label{sec:ds:float}
The paper's second deployment mode keeps the activations in float. Multiplying a float by \(\pm2^p\) adds \(p\) to its
8-bit exponent field and flips its sign bit; a zero or denormal input gives a signed zero. In the kernel this replaces
the float multiply (three DSP48E1 in the float32 builds) by a small integer adder, so a MAC is one float add, about two
DSPs. The result is exactly \(a\cdot w\) for every normal input that does not overflow (\(|a| < 2^{120}\) for 4-bit
codes), so the card's sums are checked bit for bit against float sums in the kernel's \(k\) order.

\begin{itemize}
  \item PE 32\(\times\)16, tile 512\(\times\)512, 200\,MHz: timing met, 202.5k LUT, 1063 DSP, 427.5 BRAM tiles;
        179.4, 187.6 and \textbf{189.1\,GOP/s} at \(1024^3\), \(4096^3\) and \(8192^3\) (92.3\,\% of 204.8),
        all exact. The best float32 build (multiply-add, PE 32\(\times\)8, 225\,MHz, 1320 DSP) reached
        110.9\,GFLOP/s: DenseShift weights give 1.7\(\times\) the throughput with fewer DSPs.
  \item PE 48\(\times\)16 (768 MACs, 1793 DSP) synthesised but needed 303\,086 LUTs of 298\,600: with the
        multiplier gone, the float adders' LUTs set the limit.
\end{itemize}

\subsection{Larger arrays}
\label{sec:ds:scale}
Without multipliers, the C buffers and the routing set the size. Three builds tested the limits (int8 activations,
3-bit weights, 200\,MHz unless stated):
\begin{itemize}
  \item \textbf{PE 32\(\times\)64 with int16 sums} (2048 MACs; the multiply MAC would need 2048 DSPs, more than the
        device's 1920): timing met, 190.3k LUT, 638 BRAM tiles, 40 DSP; \textbf{750.9\,GOP/s} at \(8192^3\)
        (91.7\,\%), exact. Sums wrap at 16 bits, so this suits layers whose sums fit, or a host that splits \(K\).
  \item \textbf{PE 48\(\times\)32 with int32 sums} (tile 384\(\times\)2048): 91\,\% of the BRAM tiles and
        76\,\% of the LUTs; the spread-out placement failed timing by 2.5\,ns (a PE's A register to its shift lanes,
        95\,\% routing), and the SoC and PCIe clocks failed too.
  \item \textbf{PE 32\(\times\)32 at 250\,MHz}: failed by 0.93\,ns in the HLS dataflow control
        (\code{ap\_sync}/\code{ap\_ready} fan-out), the limit already met by the large int4 arrays
        (\cref{sec:formats:wide}), not in the MAC.
\end{itemize}
With int32 sums, PE 32\(\times\)32 at 200\,MHz is therefore the practical DenseShift array on this device; its
value is the freed DSPs and smaller logic rather than a larger array. Narrower C buffers than the C bus (for example
24-bit sums) would be the next step for larger int32-class arrays.
